%% ma: L10-B14-0040
%% lop: 10
%% chuong: 5
%% bai: 14
%% dang: 10.14.2
%% loai: TLN
%% muc: VD
%% dap_an: "0,8"
%% nguon: "chuyen-de-cac-so-dac-trung-cua-mau-so-lieu-khong-ghep-nhom-toan-10-chuong-trinh-moi.pdf (GV. Phan Nhật Linh) (P04), Bài 03 Dạng 1 Phần III Câu 7"
%% kien_thuc: "Lớp 10 Bài 13 (tứ phân vị), Bài 14 (khoảng tứ phân vị)"
%% trang_thai: chinh-thuc
%% ghi_chu: "Đáp số gốc sai (0,915 do chia mỗi nửa thành 4 giá trị); đúng là ΔQ = 0,765 ≈ 0,8. Tự duyệt theo yêu cầu giáo viên 22/10/2026."
\begin{ex}
Mẫu số liệu sau đây cho biết cân nặng của $10$ trẻ sơ sinh (đơn vị: kg):
\begin{center}
\begin{tabular}{|c|c|c|c|c|}\hline
$2{,}977$ & $3{,}155$ & $3{,}920$ & $3{,}412$ & $4{,}236$\\\hline
$2{,}593$ & $3{,}270$ & $3{,}813$ & $4{,}042$ & $3{,}387$\\\hline
\end{tabular}
\end{center}
Tìm khoảng tứ phân vị của mẫu số liệu (kết quả làm tròn đến hàng phần mười).
\shortans{0,8}
\loigiai{Sắp xếp: $2{,}593$; $2{,}977$; $3{,}155$; $3{,}270$; $3{,}387$; $3{,}412$; $3{,}813$; $3{,}920$; $4{,}042$; $4{,}236$.

Nửa dưới gồm $5$ giá trị đầu nên $Q_1=3{,}155$; nửa trên gồm $5$ giá trị cuối nên $Q_3=3{,}920$.

$\Delta_Q=3{,}920-3{,}155=0{,}765\approx0{,}8$.}
\end{ex}
